% !TEX root = /home/eb/.vscode/extensions/qftrules.programme-colle/tmp/Exercice.tex
\begin{exo}[PC][1][oral]{Transmission d’une onde acoustique}
  On considérait une paroi qui séparait un milieu intérieur (air) et un milieu extérieur (air), et on pose $v=$ vo exp j $(w t-k x)$
On la modélise par une tranche infiniment fine de section $S$.
\begin{quest}
\item Donner la relation entre $S$, la masse volumique de la tranche et sa masse surfacique
\item Quelle doit être la relation entre $L$ et lambda pour pouvoir considérer que tout "vibre en bloc"? Quelle relation vérifient donc les vitesses de part et d'autre de la paroi?
\item En appliquant le PFD à une tranche de la paroi de section $S$, trouver une relation entre les pressions de part et d'autre de la paroi.
\item Trouver l'expression du coefficient de transmission en amplitude associé à la surpression, puis au coefficient de transmission énergétique $\mathrm{T}$ après avoir rappelé sa définition.
\item Tracer $\mathrm{Tdb}=\log (\mathrm{w})$, de quel type de flitre s'agit il?
\end{quest}
\end{exo}
%%%%%%%%%%%%%%%%%%%%%%%%%%%%%%%%%%%

%%%%%%%%%%%%%%%%%%%%%%%%%%%%%%%%%%%
\begin{exo}[PC][1][oral]{Dissolution d’un grain de sucre}
  On impose les vibrations de fréquence $\SI{100}{Hz}$ à une cordelette de longueur $\SI{10}{cm}$ attachée à un générateur à sa gauche et lestée par une sphère en acier à sa droite.

  On immerge ensuite la sphère dans de l’eau et on observe les vibrations du schéma ...

  Estimer le rayon de la sphère.

  \sol{

  La fréquence est celle du mode $2$, soit $f = 2f_0$ avec $f_0 = \frac{\pi c}{L}$ et $c=\sqrt{T/\mu}$. Si on passe au mode $4$ sans modifier la fréquence, c’est que $f_0$ est divisé par $2$, donc $c$ est divisé par $2$, donc $T$ est divisé par $\sqrt{2}$.

  La tension initiale vaut $T=mg$. La tension vaut ensuite $T’= mg - 4\pi \eta R v$ dans l’hypothèse où l’eau est un fluide newtonien incompressible de viscosité dynamique $\eta$. En ordre de grandeur, $v \simeq x_0 \omega = 2\pi h f$ avec $f$ la fréquence et $h$ l’amplitude des oscillations de la sphère. Donc
  $$
  T’ = \frac{T}{\sqrt{2}} = T - 8\pi^2 \eta R h f \Rar R = \frac{\sqrt{2}-1}{\sqrt{2}} \frac{mg}{8\pi^2 \eta h f}
  $$
  Masse $m=\rho \frac{4}{3}\pi R^3$ et donc
  $$
  R = \(\frac{\sqrt{2}-1}{\sqrt{2}} \frac{\rho g}{6\pi \eta h f} \)^{-1/2} \sim 8 \sqrt{\frac{\eta h f}{\rho g}}.
  $$
  On peut estimer $\eta \sim \SI{1e-3}{Pl}$, $g \sim \SI{\item1}{m.s^{-2}}$, $\rho \sim \SI{8e3}{kg.m^{-3}}$ pour l’acier, $h \simeq 4\times(10-\sqrt{10^2-2^2}) \sim \frac{4\times 2^2}{2\times 10^2} \simeq \SI{\item08}{cm} \sim \SI{1}{mm}$.
  }
\end{exo}
%%%%%%%%%%%%%%%%%%%%%%%%%%%%%%%%%%%
